\documentclass[UTF8,11pt]{ctexart}
\usepackage[a4paper,margin=24mm]{geometry}
\usepackage{amsmath,amssymb,amsthm,mathtools}
\usepackage[all]{xy}
\usepackage{xurl,hyperref,enumitem}
\hypersetup{hidelinks,breaklinks=true}
\setlength{\emergencystretch}{3em}
\newtheorem*{theorem}{Theorem}
\newtheorem*{lemma}{Lemma}
\newtheorem*{proposition}{Proposition}
\newtheorem*{corollary}{Corollary}
\theoremstyle{definition}
\newtheorem*{definition}{Definition}
\newtheorem*{remark}{Remark}
\DeclareMathOperator{\Hom}{Hom}
\DeclareMathOperator{\Ext}{Ext}
\DeclareMathOperator{\Tor}{Tor}
\DeclareMathOperator{\Spec}{Spec}
\DeclareMathOperator{\supp}{supp}
\DeclareMathOperator{\im}{im}
\DeclareMathOperator{\id}{id}
\DeclareMathOperator{\Tr}{Tr}
\DeclareMathOperator{\rank}{rank}
\newcommand{\R}{\mathbb R}
\newcommand{\C}{\mathbb C}
\newcommand{\Z}{\mathbb Z}
\newcommand{\N}{\mathbb N}
\newcommand{\Q}{\mathbb Q}
\begin{document}
\section*{023\quad 无限Galois扩张的闭子群对应}
\subsection*{来源与收录范围}
J. S. Milne, \emph{Fields and Galois Theory}, v5.00, June 2021。主目标 Theorem 7.12，印刷第95--96页；完整收入 Proposition 7.11 的闭包论证和 Propositions 7.6--7.9 的 Krull 拓扑、限制映射及紧性证明，印刷第93--95页。取得日期：2026-09-28。
\par\url{https://www.jmilne.org/math/CourseNotes/FT500.pdf}
\subsection*{作者命题与人类证明}
\paragraph{Krull topology (source, pp. 93--94).}
Let $\Omega$ be a Galois extension of $F$, and let $G=\operatorname{Aut}(\Omega/F)$. For any finite subset $S$ of $\Omega$, let
\[
G(S)=\{\sigma\in G\mid \sigma s=s\text{ for all }s\in S\}.
\]
\begin{proposition}[7.6]
There is a unique structure of a topological group on $G$ for which the sets $G(S)$ form an open neighbourhood base of $1$. For this topology, the sets $G(S)$ with $S$ $G$-stable form a neighbourhood base of $1$ consisting of open normal subgroups.
\end{proposition}
\begin{proof}
We show that the collection of sets $G(S)$ satisfies (a,b,c,d) of (7.2). It satisfies (a) because $G(S_1)\cap G(S_2)=G(S_1\cup S_2)$. It satisfies (b) and (c) because each set $G(S)$ is a group. Let $S$ be a finite subset of $\Omega$. Then $F(S)$ is a finite extension of $F$, and so there are only finitely many $F$-homomorphisms $F(S)\to\Omega$. Since $\sigma S=\tau S$ if $\sigma|_{F(S)}=\tau|_{F(S)}$, this shows that $\overline S=\bigcup_{\sigma\in G}\sigma S$ is finite. Now $\sigma\overline S=\overline S$ for all $\sigma\in G$, and it follows that $G(\overline S)$ is normal in $G$. Therefore, $\sigma G(\overline S)\sigma^{-1}=G(\overline S)\subset G(S)$, which proves (d). It also proves the second statement.
\end{proof}
The topology on $\operatorname{Aut}(\Omega/F)$ defined in the proposition is called the Krull topology. We write $\operatorname{Gal}(\Omega/F)$ for $\operatorname{Aut}(\Omega/F)$ endowed with the Krull topology, and call it the Galois group of $\Omega/F$.

If $S$ is a finite set stable under $G$, then $F(S)$ is a finite extension of $F$ stable under $G$, and hence Galois over $F$ (7.5). Therefore,
\[
\{\operatorname{Gal}(\Omega/E)\mid E\text{ finite and Galois over }F\}
\]
is a neighbourhood base of $1$ consisting of open normal subgroups.
\begin{proposition}[7.7]
Let $\Omega$ be Galois over $F$. For every intermediate field $E$ finite and Galois over $F$, the map
\[
\sigma\longmapsto\sigma|_E:\operatorname{Gal}(\Omega/F)\longrightarrow\operatorname{Gal}(E/F)
\]
is a continuous surjection (discrete topology on $\operatorname{Gal}(E/F)$).
\end{proposition}
\begin{proof}
Let $\sigma\in\operatorname{Gal}(E/F)$, and regard it as an $F$-homomorphism $E\to\Omega$. Then $\sigma$ extends to an $F$-isomorphism $\Omega\to\Omega$ (see 7.4), which shows that the map is surjective. For every finite set $S$ of generators of $E$ over $F$, $\operatorname{Gal}(\Omega/E)=G(S)$, which shows that the inverse image of $1_{\operatorname{Gal}(E/F)}$ is open in $G$. By homogeneity, the same is true for every element of $\operatorname{Gal}(E/F)$.
\end{proof}
\begin{proposition}[7.8]
The Galois group $G$ of a Galois extension $\Omega/F$ is compact and totally disconnected.
\end{proposition}
\begin{proof}
We first show that $G$ is Hausdorff. If $\sigma\ne\tau$, then $\sigma^{-1}\tau\ne1_G$, and so it moves some element of $\Omega$, i.e., there exists an $a\in\Omega$ such that $\sigma(a)\ne\tau(a)$. For any $S$ containing $a$, $\sigma G(S)$ and $\tau G(S)$ are disjoint because their elements act differently on $a$. Hence they are disjoint open subsets of $G$ containing $\sigma$ and $\tau$ respectively.

We next show that $G$ is compact. As we noted above, if $S$ is a finite set stable under $G$, then $G(S)$ is a normal subgroup of $G$, and it has finite index because it is the kernel of $G\to\operatorname{Sym}(S)$. Since every finite set is contained in a stable finite set, the argument in the last paragraph shows that the map
\[
G\longrightarrow\prod_{S\text{ finite stable under }G}G/G(S)
\]
is injective. When we endow $\prod G/G(S)$ with the product topology, the induced topology on $G$ is that for which the $G(S)$ form an open neighbourhood base of $e$, i.e., it is the Krull topology. According to the Tychonoff theorem, $\prod G/G(S)$ is compact, and so it remains to show that $G$ is closed in the product. For each $S_1\subset S_2$, there are two continuous maps $\prod G/G(S)\to G/G(S_1)$, namely, the projection onto $G/G(S_1)$ and the projection onto $G/G(S_2)$ followed by the quotient map $G/G(S_2)\to G/G(S_1)$. Let $E(S_1,S_2)$ be the closed subset of $\prod G/G(S)$ on which the two maps agree. Then $\bigcap_{S_1\subset S_2}E(S_1,S_2)$ is closed, and equals the image of $G$.

Finally, for each finite set $S$ stable under $G$, $G(S)$ is a subgroup that is open and hence closed. Since $\bigcap G(S)=\{1_G\}$, this shows that the connected component of $G$ containing $1_G$ is just $\{1_G\}$. By homogeneity, a similar statement is true for every element of $G$.
\end{proof}
\begin{proposition}[7.9]
For every Galois extension $\Omega/F$, $\Omega^{\operatorname{Gal}(\Omega/F)}=F$.
\end{proposition}
\begin{proof}
Every element of $\Omega\setminus F$ lies in a finite Galois extension of $F$, and so this follows from the surjectivity in Proposition 7.7.
\end{proof}

\begin{proposition}[7.11]
Let $\Omega$ be Galois over $F$, with Galois group $G$.
\begin{enumerate}[label=(\alph*)]
\item Let $M$ be a subfield of $\Omega$ containing $F$. Then $\Omega$ is Galois over $M$, the Galois group $\operatorname{Gal}(\Omega/M)$ is closed in $G$, and $\Omega^{\operatorname{Gal}(\Omega/M)}=M$.
\item For every subgroup $H$ of $G$, $\operatorname{Gal}(\Omega/\Omega^H)$ is the closure of $H$.
\end{enumerate}
\end{proposition}
\begin{proof}
(a) The first assertion was proved in (7.3). For each finite subset $S\subset M$, $G(S)$ is an open subgroup of $G$, and hence it is closed. But $\operatorname{Gal}(\Omega/M)=\bigcap_{S\subset M}G(S)$, and so it also is closed. The final statement now follows from (7.9).

(b) Since $\operatorname{Gal}(\Omega/\Omega^H)$ contains $H$ and is closed, it certainly contains the closure $\overline H$ of $H$. On the other hand, let $\sigma\in G\setminus\overline H$; we have to show that $\sigma$ moves some element of $\Omega^H$. Because $\sigma$ is not in the closure of $H$,
\[
\sigma\operatorname{Gal}(\Omega/E)\cap H=\varnothing
\]
for some finite Galois extension $E$ of $F$ in $\Omega$ (because the sets $\operatorname{Gal}(\Omega/E)$ form a neighbourhood base of $1$; see above). Let $\psi$ denote the surjective map $\operatorname{Gal}(\Omega/F)\to\operatorname{Gal}(E/F)$. Then $\sigma|_E\notin\psi H$, and so $\sigma$ moves some element of $E^{\psi H}\subset\Omega^H$ (apply 3.11b).
\end{proof}
\begin{theorem}[7.12]
Let $\Omega$ be Galois over $F$ with Galois group $G$. The maps
\[
H\longmapsto\Omega^H,\qquad M\longmapsto\operatorname{Gal}(\Omega/M)
\]
are inverse bijections between the set of closed subgroups of $G$ and the set of intermediate fields between $\Omega$ and $F$. Moreover,
\begin{enumerate}[label=(\alph*)]
\item the correspondence is inclusion-reversing: $H_1\supset H_2\Longleftrightarrow\Omega^{H_1}\subset\Omega^{H_2}$;
\item a closed subgroup $H$ of $G$ is open if and only if $\Omega^H$ has finite degree over $F$, in which case $(G:H)=[\Omega^H:F]$;
\item $\sigma H\sigma^{-1}\leftrightarrow\sigma M$, i.e., $\Omega^{\sigma H\sigma^{-1}}=\sigma(\Omega^H)$; $\operatorname{Gal}(\Omega/\sigma M)=\sigma\operatorname{Gal}(\Omega/M)\sigma^{-1}$;
\item a closed subgroup $H$ of $G$ is normal if and only if $\Omega^H$ is Galois over $F$, in which case $\operatorname{Gal}(\Omega^H/F)\simeq G/H$.
\end{enumerate}
\end{theorem}
\begin{proof}
For the first statement, we have to show that $H\mapsto\Omega^H$ and $M\mapsto\operatorname{Gal}(\Omega/M)$ are inverse maps.

Let $H$ be a closed subgroup of $G$. Then $\Omega$ is Galois over $\Omega^H$ and $\operatorname{Gal}(\Omega/\Omega^H)=H$ (see 7.11).

Let $M$ be an intermediate field. Then $\operatorname{Gal}(\Omega/M)$ is a closed subgroup of $G$ and $\Omega^{\operatorname{Gal}(\Omega/M)}=M$ (see 7.11).

(a) We have the obvious implications:
\[
H_1\supset H_2\Longrightarrow\Omega^{H_1}\subset\Omega^{H_2}
\Longrightarrow\operatorname{Gal}(\Omega/\Omega^{H_1})\supset\operatorname{Gal}(\Omega/\Omega^{H_2}).
\]
But $\operatorname{Gal}(\Omega/\Omega^{H_i})=H_i$ (see 7.11).

(b) As we noted earlier, a closed subgroup of finite index in a topological group is always open. Because $G$ is compact, conversely an open subgroup of $G$ is always of finite index. Let $H$ be such a subgroup. The map $\sigma\mapsto\sigma|_{\Omega^H}$ defines a bijection
\[
G/H\longrightarrow\Hom_F(\Omega^H,\Omega)
\]
(apply 7.4) from which the statement follows.

(c) For $\tau\in G$ and $\alpha\in\Omega$, $\tau\alpha=\alpha\Longleftrightarrow\sigma\tau\sigma^{-1}(\sigma\alpha)=\sigma\alpha$. Therefore, $\operatorname{Gal}(\Omega/\sigma M)=\sigma\operatorname{Gal}(\Omega/M)\sigma^{-1}$, and so $\sigma\operatorname{Gal}(\Omega/M)\sigma^{-1}\leftrightarrow\sigma M$.

(d) Let $H\leftrightarrow M$. It follows from (c) that $H$ is normal if and only if $M$ is stable under the action of $G$. But $M$ is stable under the action of $G$ if and only it is a union of finite extensions of $F$ stable under $G$, i.e., of finite Galois extensions of $G$. We have already observed that an extension is Galois if and only if it is a union of finite Galois extensions.
\end{proof}

\subsection*{整理说明（非作者正文）}
只计无限 Galois 对应一个问题，所附拓扑与闭包性质是其核心证明，不另计数。标准先修：有限 Galois 理论（包括不动域定理）、代数可分扩张嵌入延拓（7.4）、正规可分扩张是有限 Galois 扩张的并（7.3--7.5）、Tychonoff 定理和拓扑群的邻域基准则（7.2）。7.2 所用条件为有限交可细化、乘积可细化、逆可细化及共轭可细化；这些是拓扑群的基本构造条件，而非本题的闭子群对应。

难度在于有限不动域理论如何通过 Krull 邻域、有限商的乘积和闭包识别扩展到无限情形；没有把这一对应本身列作先修。所附作者主证明 (d) 最后写 ``finite Galois extensions of $G$''，依前句底域为 $F$；保留原文并说明该符号问题，不暗中以助手证明替换。群元素和限制映射分别统一排作 $\tau$、$\psi$，仅作记号排版。已取得完整文字；第96页截图接口失败，未声称逐页图像核对，未保存下载 PDF 原件。
\subsection*{可修改的简短标注（非作者正文）}
$c$：无限代数扩张的中间域分类

$a$：不动域；Krull 拓扑；有限商群；Tychonoff 紧性；闭包与有限邻域分离；嵌入延拓

\texttt{cross\_theory\_status = yes}。把中间域识别转成自同构子群的拓扑闭包问题；有限商提供可分离的邻域，紧性又把开子群与有限指数相联系，返回有限子扩张的次数与整体域对应。位置：Propositions 7.8、7.11 与 Theorem 7.12。

全部标注均为非穷尽、可修改初稿；未标注不表示未使用。
\subsection*{许可与修改记录}
Copyright 1996--2021 J. S. Milne。v5.00 按 CC BY-NC-SA 4.0 发布，本文摘录及整理沿用同一许可。版本与许可来源见 \texttt{sources/milne\_ft500/SOURCE\_AND\_LICENSE.md}。2026-09-28：助手转录、独立排版及初稿标注；不暗示原作者认可。\url{https://creativecommons.org/licenses/by-nc-sa/4.0/}。
\end{document}
